\Article{30}
Der Satz des vorigen Artikels bedarf in einem besonderen Fall, den wir nicht stillschweigend übergehen dürfen, einer gewissen Abänderung. Wir haben nämlich stillschweigend vorausgesetzt, dass die Oberfläche des Gefäßes längs der ganzen Grenze $P$ stetig gekrümmt ist, sodass es in jedem Punkt dieser Grenze \textit{genau eine} Tangentialebene an die Gefäßoberfläche gibt. Wird die Stetigkeit der Krümmung an einem einzelnen Punkt der Linie $P$ unterbrochen, sei es durch eine Spitze an diesem Punkt oder durch eine Kante, die die Linie $P$ kreuzt, so sieht man leicht, dass unsere Folgerung dadurch unverändert bleibt. Anders verhält es sich, wenn die Stetigkeit der Krümmung auf einem endlichen Teil der Linie $P$ unterbrochen ist, also wenn die Gefäßoberfläche auf einem endlichen Teil der Linie $P$ oder sogar längs dieser ganzen Linie eine Kante aufweist. In jedem Punkt eines solchen Teils gibt es dann zwei Ebenen,
\SourceLeaf{74}die die Gefäßoberfläche berühren. Die eine gehört zum freien Teil der Gefäßoberfläche, die andere zum Teil $T$. Behalten wir also das Zeichen $i$ für den Winkel zwischen der ersten Ebene und der Tangentialebene an die Oberfläche $U$ bei und bezeichnen mit $k$ den Winkel zwischen dieser Tangentialebene und der zweiten Ebene, so gilt nicht mehr $i+k=180^0$. Die Summe ist größer oder kleiner, je nachdem, ob die Kante konvex oder konkav ist. Das Element der Variation $\delta W$ wird bei einer virtuellen Bewegung, bei der die Flüssigkeit über die Grenze $P$ hinausfließt, weiterhin ausgedrückt durch
\[
\alpha\alpha\mathrm{d}P.\delta e.\sin(5,6).(\cos A-\cos i)
\]
bei einer virtuellen Bewegung, bei der sich die Flüssigkeit diesseits der Grenze $P$ zurückzieht, ist das Element dieser Variation nun
\[
-\alpha\alpha\mathrm{d}P.\delta e.\sin(5,6).(\cos A+\cos k)
\]
Damit der Wert von $W$ keine negative Variation erfahren kann, darf der Wert von $\cos A-\cos i$ nicht negativ und der von $\cos A+\cos k$ nicht positiv sein. Es muss also gelten
\[
\begin{aligned}
&\text{entweder}\quad i=A,\qquad\text{entweder}\quad i>A\\
&\text{und entweder}\quad k=180^0-A,\qquad\text{entweder}\quad k>180^0-A
\end{aligned}
\]
Im Gleichgewichtszustand kann daher $i+k<180^0$ nicht eintreten. Gleichbedeutend damit ist, dass \textit{im Gleichgewichtszustand die Grenze der freien Flüssigkeitsoberfläche $U$ nicht über eine endliche Strecke auf einer konkaven Kante der Gefäßoberfläche liegen kann.} Fällt dagegen ein Teil dieser Grenze mit einer konvexen Kante zusammen, so ist für das Gleichgewicht notwendig und hinreichend, dass der Winkel zwischen den Tangentialebenen an die Flüssigkeit und an das Gefäß zwischen $A$ und $A+a$ (einschließlich) liegt, wenn er außerhalb der Flüssigkeit gemessen wird, oder zwischen $180^0-A$ und $180^0-A+a$, wenn er innerhalb der Flüssigkeit gemessen wird. Dabei bezeichnen wir mit $180^0-a$ allgemein den Winkel zwischen den beiden Ebenen, die die Gefäßoberfläche beiderseits der Kante im jeweiligen Punkt berühren, soweit dieser Winkel auf der Gefäßseite gemessen wird.

\Article{31}
Die Konstanten $\alpha\alpha$, $\beta\beta$, deren Verhältnis den Winkel $A$ bestimmt, hängen von den Funktionen $f,F$ ab und lassen sich in gewissem Sinne als Maße für die Stärke der Molekularkräfte auffassen, die die Teilchen der Flüssigkeit und des Gefäßes ausüben. Sind diese Funktionen so beschaffen, dass $fx,Fx$ in einem bestimmten, vom Abstand $x$ unabhängigen Verhältnis stehen, etwa wie $n$ zu $N$, so können wir offenbar setzen
\SourceLeaf{75}$\alpha\alpha:\beta\beta=cn:CN$. Das heißt, die Konstanten $\alpha\alpha$, $\beta\beta$ sind proportional zu den Anziehungskräften, die im gleichen Abstand zwei Moleküle gleichen Volumens ausüben, das eine aus der Flüssigkeit, das andere aus dem Gefäß. Nun wird der Winkel $A$ spitz, gleich einem rechten Winkel, stumpf oder gleich zwei rechten Winkeln, je nachdem, ob $\beta\beta<\frac12\alpha\alpha$, $\beta\beta=\frac12\alpha\alpha$, $\beta\beta>\frac12\alpha\alpha$ aber $<\alpha\alpha$ oder $\beta\beta=\alpha\alpha$ gilt. Unter dieser Annahme, die zwar unbegründet sein mag, aber dem Wahrscheinlichen nicht widerspricht, müssen wir sagen, dass der erste Fall eintritt, wenn die gegenseitige Anziehung der Flüssigkeitsteile größer ist als das Doppelte der Anziehung der Gefäßteile auf die Flüssigkeit. Der zweite Fall tritt ein, wenn die erste Anziehung doppelt so groß ist wie die zweite. Der dritte tritt ein, wenn die erste Anziehung zwar größer als die zweite, aber kleiner als deren Doppeltes ist. Der vierte tritt schließlich ein, wenn beide Anziehungen gleich sind. Ein Beispiel für den ersten Fall ist Quecksilber in Glasgefäßen.

